GRAIL predicts xenobiotic metabolite structures in three stages. A learned generator scores every
rule in a bank of $\nbank$ SMIRKS templates against the substrate, estimating the probability
$P(r\mid s)$ that rule $r$ applies to substrate $s$. RDKit then applies each selected rule and
enumerates the products, recording which rule produced which candidate. A structural filter scores
each substrate--product pair, and the deployed system ranks candidates by the product of the
filter and generator scores, combining a judgement about the plausibility of the pair with the
generator's confidence in the rule that made it.

Both learned stages train on positive-unlabelled data. Annotated metabolites are the only
positives; rule-applicable products that carry no annotation are unlabelled, not negative,
because a missing annotation does not establish that a transformation does not occur. The filter
exposes its raw classifier outputs so that a non-negative positive-unlabelled surrogate is never
applied on top of a sigmoid when that estimator is selected, and the generator down-weights
applicable but unobserved rules instead of penalising them as hard negatives.

Both stages share a featurisation: molecular graphs with 16-dimensional nodes and 18-dimensional
edges, and a 1024-bit Morgan fingerprint per molecule. Unless stated otherwise, matching between
predicted and reference structures uses the tautomer-aware InChIKey convention throughout this
paper.

The three stages are separately inspectable: the selected rule, the enumerated product, and the
filter's judgement of the pair. What the architecture contributes is interpretable rule selection
paired with a positive-unlabelled filter, not recall above other metabolite predictors.

\subsection{The generator}

The generator scores the entire bank in one forward pass. A GATv2 encoder with three
message-passing layers, node dimensions $16\to192\to256\to128$, batch normalisation, ReLU and
dropout of 0.1 at each layer, and a mean-pooled readout, embeds the substrate graph into 128
dimensions. A parallel multilayer perceptron embeds its Morgan fingerprint, and a gated fusion
combines the two into the substrate representation. Each rule is embedded independently by a
rule-graph encoder of the same family, with hidden dimensions $128\to192\to128$, alongside a
six-dimensional branch of global rule features. The full bank is re-encoded on every forward pass,
which dominates the compute cost and is why inference time scales with bank size.

The logit for a rule sums four learned terms, each scaled by its own learned scalar. A
cross-attention term lets the rule embedding query the substrate's atom states through
single-head attention over node keys and values, producing a rule-specific local context that is
concatenated with the pooled embeddings and passed through a three-layer perceptron. A global term
is the cosine similarity between the substrate and rule embeddings, and a local term the cosine
similarity between the attention context and the rule. A fourth term reflects whether the rule's
SMARTS pattern matches. To this sum the model adds a per-rule bias and a frequency-prior term at
weight 0.4, the log-odds that the rule yields a true metabolite as estimated from training
statistics, then subtracts a penalty of 7.5 for rules whose SMARTS does not match, which
effectively masks inapplicable rules before the sigmoid. Where a rule fires at several sites, the
per-firing probabilities are combined by noisy-OR.

Training is multi-label rule classification with a margin-ranking auxiliary at weight 0.25 and
margin 0.45. Rules that match the substrate but carry no annotation are down-weighted by half
rather than treated as negatives, which is what makes the objective positive-unlabelled.
Optimisation uses Adam at a learning rate of $10^{-4}$ for at most 8 epochs with early stopping.

\subsection{The filter}

The deployed filter, which is the checkpoint behind every GRAIL number reported here, embeds the
substrate and the candidate product independently. Each passes through its own three-layer GATv2
encoder over the molecule's graph, with dimensions $16\to192\to256\to128$ and a mean-pooled
readout. The two 128-dimensional embeddings are concatenated with both Morgan fingerprints, and
the resulting 2304-dimensional vector is mapped through a three-layer perceptron
($\to128\to64\to1$, ReLU, dropout 0.1) to a single logit.

An alternative variant encodes one merged substrate--product graph, whose 18-dimensional nodes
are joined by cross-edges placed according to an element-aware maximum-common-substructure
correspondence between atoms, not by index order, so that chemically meaningful atom
mappings are preserved across the reaction. We measured the two rather than assuming which should
be the default, and report the comparison below.

The estimator is selectable: with the non-negative positive-unlabelled risk it consumes logits
and with binary cross-entropy it consumes probabilities, and the released configuration trains the
filter under the second. Optimisation uses Adam at a learning rate of $10^{-4}$ for at most 8 epochs
with early stopping at a patience of 7.

\subsection{Choosing between the two filters}

Both variants were trained on an identical 800-substrate subset, with the same negatives, seed and
optimiser, differing only in which encoder they used, and evaluated on the full clean test split
with the same generator and candidate pool.

Independent encoders win on accuracy. Recall@15 is 0.366 for the independent variant against 0.357
for the merged-graph variant, a paired difference of $-0.0090$ with a 95\% paired-bootstrap interval of
$[-0.017,-0.001]$ over $n=\ntest$ substrates; at recall@5 the margin widens, 0.282 against 0.259.

They also win on cost by roughly twenty-fold. Computing a maximum common substructure for every
pair makes full-data training take about 14 hours against roughly 40 minutes, and imposes the same
computation at every candidate at inference. Accuracy and cost point the same way. The merged-graph variant is implemented and available, and is not deployed because it loses on both counts.

\subsection{How the rule bank was built}

The rule-based system GRAIL is compared against, SyGMa, uses hand-authored reaction
rules. GRAIL's bank is mined automatically from annotated substrate--metabolite pairs, with each
candidate template validated against the pair that produced it.

Mining reads only the clean, molecule-disjoint training split. For each annotated pair it first
locates a ring-aware maximum common substructure, using RDKit's FMCS with ring atoms
matched only to ring atoms, rings completed, any bond order allowed to match, and atoms compared
by element; the match must cover at least 40\% of the smaller molecule's atoms. From that common
substructure it derives a reaction centre. Three kinds of atom are flagged: matched atoms whose element,
formal charge, hydrogen count, aromaticity or degree changed; atom pairs whose connecting bond
order changed; and leaving or entering atoms together with their matched neighbours. Several
substrate--product match combinations are tried, and the one giving the smallest centre is kept.
The centre is then
expanded by one bond to capture its immediate reactive environment, and the expanded fragments are
atom-mapped by their positional correspondence in the maximum common substructure into a candidate
SMIRKS template.

Two gates follow. A template survives only if applying it to its own source substrate regenerates
the annotated product after canonicalisation. The surviving, deduplicated templates then pass a
selectivity filter that rejects templates producing more than 200 distinct products on average
across a 50-substrate sample, and templates that are both very general and very prolific,
applicable to more than 90\% of sampled substrates while still averaging more than 50
products.

A second atom-mapping route is also implemented: it detects the reaction centre with RASCAL MCES,
falling back to FMCS, expands by one bond as above, and emits a mapped SMIRKS using RXNMapper, a
neural attention-based atom mapper. Both a rule-based and a neural mapping route are therefore
available.

Because mining reads the training split alone, no test metabolite can enter a mined rule, so the
coverage ceiling reported in the main text remains a held-out quantity.

\subsection{What the bank contains}

The deployed bank of $\nbank$ templates is the deduplicated union of four earlier curated banks,
of 473, 656, 500 and 1{,}051 rules, together with the newly mined templates.

Support across mined templates is heavily skewed. The positional miner alone extracts 5{,}856
unique self-tested templates from the training positives. Of these, 4{,}274 (73\%)
rest on a single training pair, and only 457 on five or more; the single best-supported template
covers 1{,}531 pairs. The bank is therefore dominated by highly specific, near-single-example
templates, not by broad general reactions.

Re-running the miner independently over the same training split re-derives templates that are all
already in the bank: 5{,}856 of 5{,}856, none new. Mining is deterministic and, on this corpus,
saturated at its source. The coverage ceiling of $\ceiling$ is bounded by the diversity of the
training reactions, not by incomplete mining, so extending coverage requires new reaction corpora
or more general templates rather than further mining of the same data.

The selectivity filter is the construction-time counterpart of the breadth--precision trade
measured in Appendix~\ref{app:props}. Rejecting over-general templates when the bank is built and
narrowing the number of rules applied at inference act on the same quantity, one before
deployment and one at the operating point.
